\documentclass{article}
\usepackage{amsmath, amssymb, amsthm}
\usepackage{hyperref}
\usepackage{geometry}
\geometry{margin=1in}

\title{Erd\H{o}s Problem \#10}
\date{Last updated: 2025-08-31}
\author{Source: erdosproblems.com}

\begin{document}
\maketitle

\noindent\textbf{Status:} OPEN \quad \textbf{No prize}

\noindent\textbf{Tags:} number theory, additive basis, primes

\medskip

\noindent\textbf{Problem Statement:}

Is there some $k$ such that every large integer is the sum of a prime and at most $k$ powers of 2?




Erd\H{o}s described this as 'probably unattackable'. In \cite{ErGr80} Erd\H{o}s and Graham suggest that no such $k$ exists, although in \cite{Er80} Erd\H{o}s conjectured with 'trepidation' that such a $k$ does exist. Gallagher \cite{Ga75} has shown that for any $\epsilon&gt;0$ there exists $k(\epsilon)$ such that the set of integers which are the sum of a prime and at most $k(\epsilon)$ many powers of 2 has lower density at least $1-\epsilon$. Granville and Soundararajan \cite{GrSo98} have conjectured that at most $3$ powers of 2 suffice for all odd integers, and hence at most $4$ powers of $2$ suffice for all even integers. (The restriction to odd integers is important here - for example, Bogdan Grechuk has observed that $1117175146$ is not the sum of a prime and at most $3$ powers of $2$, and pointed out that parity considerations, coupled with the fact that there are many integers not the sum of a prime and $2$ powers of $2$ (see [9] ) suggest that there exist infinitely many even integers which are not the sum of a prime and at most $3$ powers of $2$). See also [9] , [11] , and [16] .




References

[Er80] Erd\H{o}s, Paul, A survey of problems in combinatorial number theory . Ann. Discrete Math. (1980), 89-115.

[ErGr80] Erd\H{o}s, P. and Graham, R., Old and new problems and results in combinatorial number theory . Monographies de L'Enseignement Mathematique (1980).

[Ga75] Gallagher, P. X., Primes and powers of 2 . Invent. Math. (1975), 125-142.

[GrSo98] Granville, A. and Soundararajan, K., A Binary Additive Problem of Erd\H{o}s and the Order of $2$ mod $p^2$ . The Ramanujan Journal (1998), 283-298.

\noindent\textbf{Related OEIS sequences:} A387053


\bigskip
\noindent\small{Source: \url{https://www.erdosproblems.com/10}}
\end{document}
